\documentclass[a4paper,12pt]{jsarticle}
\usepackage[dvipdfmx,margin=24mm]{geometry}
\usepackage[dvipdfmx]{graphicx}
\usepackage{amsmath,amssymb,amsthm}
\usepackage{enumitem}
\usepackage{tikz}
\usetikzlibrary{arrows.meta}

\setlength{\parindent}{0pt}
\setlength{\parskip}{0.35em}
\setlist[itemize]{leftmargin=2em,itemsep=0.2em,topsep=0.3em}

\newcommand{\feedbackqa}[2]{%
\item[]
\begin{tabular}[t]{@{}p{1.8em}@{\hspace{0.4em}}p{\dimexpr\linewidth-2.2em\relax}@{}}
Q. & #1 \\
A. & #2
\end{tabular}
\par\smallskip\hrulefill
}

\newcommand{\feedbacknote}[1]{%
\item[] #1\par\smallskip\hrulefill
}

\newcommand{\feedbackrule}{%
\item[] \leavevmode\hrulefill
}

\theoremstyle{definition}
\newtheorem*{definition}{定義}
\newtheorem*{theorem}{定理}

\begin{document}

\begin{center}
{\LARGE 数学1A 第4回レジュメ}\\[0.4em]
ロピタルの定理とマクローリン展開
\end{center}

\section*{フィードバック}

\begin{itemize}[leftmargin=1.5em,itemsep=0.7em,topsep=0.3em]
\feedbackqa{勉強時間の目安はどのくらいか。}{理解度にもよりますが、この授業ではまず「解き方の原理」を理解することを目標にしています。解き方を覚え、なぜその方法でうまくいくのかが分かったあとに、いくつか演習問題を解くのがよいと思います。あくまで個人的な目安ですが、授業時間と同じ1時間30分くらいでしょうか。}

\feedbackqa{コーシーの平均値の定理で，なぜ
\[
F(x)=\frac{f(b)-f(a)}{g(b)-g(a)}\bigl(g(x)-g(a)\bigr)-\bigl(f(x)-f(a)\bigr)
\]
のような関数を考えるのか。}{私自身も，この種の関数を「自然に思いつく」感覚はまだよくわかりません。いろいろ調べると説明は見つかりますが，どれも腑に落ちません。結局のところ，ロルの定理をうまく適用できるように，条件を満たす関数を人工的に作っている，と理解するのがいちばん近いように思います。}

\feedbackqa{$a<b$ のとき，なぜ
\[
\frac{b-a}{b}\leq \frac{b-a}{a}
\]
が成り立つか。}{$0<a<b$ なので $a,b$ はどちらも正です。このとき $a<b$ なら $\frac1b<\frac1a$ ですから，正の数 $b-a$ を掛ければ
\[
\frac{b-a}{b}\leq \frac{b-a}{a}
\]
が従います。分母が大きいほど，全体の値は小さくなると考えてもよいです。}

\feedbackqa{小テストで $x>e$ という条件はどこで使っていたのですか。}{鋭い指摘です。実は後で気づきましたが，その条件はいらないです、、}

\feedbackqa{コーシーの平均値の定理と，高校で学ぶ平均値の定理は何が違うか。}{高校で学ぶ平均値の定理は，コーシーの平均値の定理の特別な場合です。コーシーの平均値の定理では2つの関数 $f,g$ を扱えるので，より一般的な形になっています。ロピタルの定理やテイラーの定理の証明で重要な役割を果たします。}

\feedbackqa{ロルの定理とコーシーの平均値の定理の違いは何ですか。}{ロルの定理は，1つの関数について $f(a)=f(b)$ のときに，途中で $f'(c)=0$ となる点があることを述べる定理です。これに対してコーシーの平均値の定理は，2つの関数 $f,g$ を用いて
\[
\frac{f(b)-f(a)}{g(b)-g(a)}=\frac{f'(c)}{g'(c)}
\]
という関係を与える定理です。}

\feedbackqa{コーシーの平均値の定理は何に使うか。}{今回扱ったロピタルの定理の証明に使いますし，特に重要なのはテイラーの定理の証明です。今後の解析でもたびたび登場します。}

\feedbackqa{証明を書くときは，何から手をつければよいか。}{特別な決まりがあるわけではありませんが，まずは「条件から何が言えるか」を丁寧に整理するのが大切です。条件から導ける事実を少しずつ積み上げていくと，目標にたどり着けることが多いです。}

\feedbackqa{対数微分では，$f(x)\neq 0$ と毎回書かなくてもよいか。}{できるだけ書いたほうがよいと思います。授業では省略してしまうこともありますが，答案では条件を明記するほうが丁寧です。ただし，文脈上明らかな場合には，省略しても大きな問題にならないこともあります。}

\feedbackqa{証明の最後に書くニコチャンマークは，テストでも使ってよいか。}{この授業では構いません。ただし，他の授業でも同じように許されるかどうかは担当教員によります。}

\feedbackqa{平均値の定理と中間値の定理は名前が似ていて紛らわしい。}{たしかに紛らわしいですね。}

\feedbackqa{$\log|x^2+1|$ の絶対値は，自分で外してよいか。}{絶対値の中身が明らかに非負であると分かるなら，自分で外して構いません。たとえば $x^2+1>0$ は明らかです。一方で，$\log(1+\cos x)$ のように符号がすぐには判断しにくい場合は，一言説明を添えるとより丁寧です。}
 
\end{itemize}
\section*{小テスト}

\subsection*{問題}

\begin{enumerate}
\item （50点）対数微分を用いて，次の関数を微分せよ。
\[
f(x)=\frac{x}{(x^2+1)^4(2x^4+1)^8}
\]

\item （50点）$0<a<b$ に対して，次を示せ。ただし，コーシーの平均値の定理を用いること。
\[
\frac{b-a}{b}\leq \log b-\log a\leq \frac{b-a}{a}
\]
\end{enumerate}

\subsection*{解答}

\begin{enumerate}
\item $x\neq 0$ として対数をとると
\[
\log |f(x)|=\log |x|-4\log(x^2+1)-8\log(2x^4+1)
\]
である。したがって
\[
\frac{f'(x)}{f(x)}
=\frac{1}{x}-\frac{8x}{x^2+1}-\frac{64x^3}{2x^4+1}.
\]
よって
\[
f'(x)=
\frac{x}{(x^2+1)^4(2x^4+1)^8}
\left(
\frac{1}{x}-\frac{8x}{x^2+1}-\frac{64x^3}{2x^4+1}
\right).
\]
また $x=0$ では，定義により
\[
f'(0)=\lim_{h\to 0}\frac{f(h)-f(0)}{h}
=\lim_{h\to 0}\frac{1}{(h^2+1)^4(2h^4+1)^8}
=1
\]
である。

\item 関数 $f(x)=\log x$, $g(x)=x$ に，区間 $[a,b]$ でコーシーの平均値の定理を用いる。ある $c\in(a,b)$ が存在して
\[
\frac{\log b-\log a}{b-a}
=\frac{f'(c)}{g'(c)}
=\frac{1/c}{1}
=\frac{1}{c}
\]
となる。$a<c<b$ より
\[
\frac{1}{b}<\frac{1}{c}<\frac{1}{a}
\]
である。$b-a>0$ を掛けると
\[
\frac{b-a}{b}\leq \log b-\log a\leq \frac{b-a}{a}
\]
を得る。
\end{enumerate}

\section{ロピタルの定理}

まず，コーシーの平均値の定理を思い出す。

\begin{theorem}[コーシーの平均値の定理]
$f,g$ が閉区間 $[a,b]$ で連続，開区間 $(a,b)$ で微分可能とする。このとき，ある $c\in(a,b)$ が存在して
\[
\frac{f(b)-f(a)}{g(b)-g(a)}=\frac{f'(c)}{g'(c)}
\]
が成り立つ。ただし $g'(c)\neq 0$ とする。
\end{theorem}

これを用いると，$0/0$ 型の極限に対するロピタルの定理が得られる。

\begin{theorem}[ロピタルの定理]
$a$ の近くで $f,g$ が微分可能で，$f(a)=g(a)=0$ とする。さらに $g'(x)\neq 0$ が $a$ の近くで成り立ち，
\[
\lim_{x\to a}\frac{f'(x)}{g'(x)}=L
\]
が存在するとする。このとき
\[
\lim_{x\to a}\frac{f(x)}{g(x)}=L
\]
が成り立つ。
\end{theorem}

\begin{proof}
$x\neq a$ を $a$ の十分近くの点とする。$[a,x]$ または $[x,a]$ にコーシーの平均値の定理を適用すると，ある $c$ が $a$ と $x$ の間に存在して
\[
\frac{f(x)-f(a)}{g(x)-g(a)}=\frac{f'(c)}{g'(c)}
\]
となる。ここで $f(a)=g(a)=0$ より
\[
\frac{f(x)}{g(x)}=\frac{f'(c)}{g'(c)}
\]
である。$x\to a$ とすると $c\to a$ だから，右辺は仮定より $L$ に収束する。したがって
\[
\lim_{x\to a}\frac{f(x)}{g(x)}=L
\]
を得る。
\end{proof}

\paragraph{応用1.}
\[
\lim_{x\to 0}\frac{e^x-1}{x}
=\lim_{x\to 0}\frac{e^x}{1}=1.
\]

\paragraph{応用2.}
\[
\lim_{x\to 0}\frac{\log(1+x)}{x}
=\lim_{x\to 0}\frac{1/(1+x)}{1}=1.
\]

\section{高階微分}

関数 $f$ に対して
\[
f''(x)=(f'(x))'
\]
と定義する。さらに $n\geq 2$ に対して
\[
f^{(n)}(x)=\left(f^{(n-1)}(x)\right)'
\]
と定義し，これを $f$ の $n$ 階微分という。

例として，
\[
f(x)=e^x
\]
のときは
\[
f'(x)=e^x,\qquad f''(x)=e^x,\qquad f^{(n)}(x)=e^x
\]
である。また，
\[
f(x)=x^2+4x+1
\]
のときは
\[
f'(x)=2x+4,\qquad f''(x)=2,\qquad f^{(n)}(x)=0\quad(n\geq 3)
\]
となる。

\section{級数}

数列 $\{a_n\}_{n=0}^{\infty}$ に対して
\[
S_n=a_0+a_1+\cdots+a_n=\sum_{k=0}^n a_k
\]
を第 $n$ 部分和という。部分和列 $\{S_n\}$ がある値 $S$ に収束するとき，
\[
\sum_{n=0}^{\infty} a_n=S
\]
と書き，この級数は収束するという。

\paragraph{例1. 幾何級数}
$|r|<1$ のとき
\[
1+r+r^2+\cdots+r^n=\frac{1-r^{n+1}}{1-r}
\]
なので，$n\to\infty$ とすると
\[
\sum_{n=0}^{\infty} r^n=\frac{1}{1-r}
\]
を得る。

\paragraph{例2. 発散の例}
\[
1-1+1-1+\cdots
\]
の部分和は $1,0,1,0,\dots$ と振動するので収束しない。

\section{マクローリン展開}

数列 $\{a_n\}$ を用いて
\[
\sum_{n=0}^{\infty} a_n x^n
\]
の形で表される式を $x=0$ を中心とするべき級数という。

たとえば $|x|<1$ では
\[
\frac{1}{1-x}=1+x+x^2+x^3+\cdots=\sum_{n=0}^{\infty} x^n
\]
が成り立つ。

では，一般の関数 $f$ が
\[
f(x)=\sum_{n=0}^{\infty} a_n x^n
\]
と表せるとすると，係数 $a_n$ はどう決まるだろうか。

両辺を順に微分して $x=0$ を代入すると
\[
f(0)=a_0,
\qquad
f'(0)=a_1,
\qquad
f''(0)=2!a_2,
\qquad
\dots,
\qquad
f^{(n)}(0)=n!a_n
\]
となる。したがって
\[
a_n=\frac{f^{(n)}(0)}{n!}
\]
であるはずだと予想される。

\begin{theorem}[マクローリン展開]
適当な条件のもとで，関数 $f$ は
\[
f(x)=\sum_{n=0}^{\infty}\frac{f^{(n)}(0)}{n!}x^n
\]
と表される。これを $f$ のマクローリン展開という。
\end{theorem}

代表例として，
\[
e^x=\sum_{n=0}^{\infty}\frac{x^n}{n!}.
\]
がある。
\end{document}
